\section{Related work}
\subsection{Language models as judges and short-answer graders}
The use of language models as judges motivates both the flexibility and the risks of this design. Zheng et al.\ study model-based evaluation of conversational responses and discuss biases in model judgements \citep{zheng2024judging}. That setting differs from exam marking: an educational rubric assigns credit to specific demonstrated ideas, and a plausible response can still contain a factual error. The present system therefore constrains model output with criterion identifiers, mark bounds and evidence checks instead of treating a generated score as self-validating.

Automated short-answer grading also predates current generative models. Mohler et al.\ use semantic similarity and dependency-graph alignment to connect student responses to reference answers \citep{mohler2011learning}. This provides a useful distinction between lexical overlap and the meaning of an answer. Auto-Assessment retains a lexical baseline as a deliberately limited reference point. Its scripted offline provider similarly tests control flow rather than semantic judgement; successful dry runs are not evidence that the language-model grader understands an answer.

\subsection{Indic OCR and handwritten documents}
The implementation integrates Bodhan's hosted IndicOCR interface, local IndicOCR weights and precomputed OCR JSON, with provider adapters converting results to common page and line contracts. The local adapter retains reading order, boxes, confidence and block kinds where available. The combination of page layout and text recognition is relevant to scripts with multiple questions, equations and mixed writing systems.
% TODO cite: verify a primary Bodhan / AI4Bharat IndicOCR publication or model card before adding a bibliographic claim.
No OCR benchmark result is imported from the archived project report. Language coverage advertised by a provider must be distinguished from accuracy on the project's own handwriting and scanning conditions. This report evaluates the latter only when paired scans and transcripts are available.

\subsection{Selective prediction}
Selective classification formalises a system's ability to abstain rather than answer every input, and examines the trade-off between coverage and risk \citep{geifman2017selective}. Auto-Assessment adopts that operational perspective: the teacher should see how accurate the automatically accepted subset is and how much work remains for manual review. Its confidence is a minimum of diagnostic signals, not a calibrated probability of correctness. A coverage curve can assess the usefulness of that heuristic, but cannot by itself establish fairness or eliminate systematic errors shared by the graders.
